\chapter{Introduction}
\label{cap:Introduction}

Numerical methods constitute a fundamental area of Applied Mathematics and
Scientific Computing, providing techniques for obtaining approximate solutions
to mathematical problems for which analytical solutions are unavailable,
unknown, or computationally impractical. Problems involving roots of nonlinear
equations, systems of equations, numerical integration, interpolation, and
differential equations are commonly encountered in engineering, physics,
computer science, and mathematical modeling
\cite{burden2016numerical,chapra2021numerical}.

Although numerical methods are defined through mathematical formulations,
their practical application depends directly on computational implementation.
Consequently, several computational aspects can influence their behavior,
including the programming language, programming paradigm, control structures,
numerical representation, compiler or interpreter, and the limitations
inherent to finite-precision arithmetic
\cite{higham2002accuracy}.

The same numerical method may therefore assume significantly different
computational structures depending on the programming paradigm adopted.
Programming paradigms provide different mechanisms for expressing computation,
state, control flow, and abstraction. In imperative programming, algorithms are
typically described as sequences of instructions that explicitly modify the
state of a program. Mutable variables, conditional statements, and iterative
control structures are therefore natural mechanisms for representing
iterative numerical procedures
\cite{scott2015programming,kernighan1988c}.

Functional programming, in contrast, emphasizes the evaluation and composition
of functions rather than the explicit modification of program state. Concepts
such as immutability, recursion, higher-order functions, and function
composition provide alternative ways of expressing computational processes.
Functional programming has also been associated with modularity and with the
ability to decompose programs through higher-order functions and lazy
evaluation
\cite{hughes1989functional}.

This distinction becomes particularly relevant when implementing iterative
numerical algorithms. Consider, for example, an iterative sequence defined by

\begin{equation}
    x_{n+1} = g(x_n).
\end{equation}

In an imperative implementation, this process can naturally be represented by
repeatedly updating a variable inside a loop. In a functional implementation,
the same recurrence can be expressed through successive function applications,
recursive definitions, or transformations over immutable values.

Although these representations describe the same mathematical recurrence,
their computational implementations can differ in terms of code structure,
expressiveness, modularity, memory behavior, and execution time. Therefore,
the choice of programming paradigm is not merely a syntactic decision; it
also determines how the mathematical algorithm is represented and executed by
the computer.

This problem is particularly relevant for iterative numerical methods, in
which the same operation may be performed hundreds, thousands, or millions of
times. In such cases, differences introduced by the programming language,
runtime system, compiler optimizations, memory-management strategy, or
evaluation model may significantly affect the final computational cost.

Another fundamental issue arises from the fact that digital computers cannot,
in general, represent arbitrary real numbers exactly. Instead, real-valued
quantities are represented using finite sets of machine numbers, most commonly
through floating-point arithmetic
\cite{higham2002accuracy,ieee7542019}.

The IEEE 754 standard defines widely adopted formats and operations for binary
and decimal floating-point arithmetic, including rules for rounding, exceptional
values, and arithmetic behavior \cite{ieee7542019}. Because the number of
available bits is finite, many real numbers cannot be represented exactly and
must instead be approximated.

For example, a real number $x$ that possesses a finite decimal representation
may have an infinite binary expansion. When represented using a finite number
of bits, its stored value can be written conceptually as

\begin{equation}
    \tilde{x} = x + \delta,
\end{equation}

where $\delta$ denotes the error introduced by finite-precision representation.

Such errors are especially important in iterative numerical algorithms because
results produced at one iteration are commonly used as inputs to subsequent
iterations. As a result, rounding errors can propagate, accumulate, or be
amplified depending on the numerical properties of the algorithm
\cite{higham2002accuracy}.

The research presented in this work is therefore situated at the intersection
of three principal topics: numerical methods, programming paradigms, and
floating-point arithmetic.

The central research question can be formulated as follows:

\begin{quote}
\textit{How do different programming paradigms and programming languages
influence the implementation, computational performance, readability, and
numerical behavior of algorithms used in numerical methods?}
\end{quote}

To investigate this question, implementations in C, Python, and Haskell are
considered. C provides a predominantly imperative programming model, offering
explicit control over variables, control structures, data representation, and
memory organization \cite{kernighan1988c}.

Python is a general-purpose programming language widely used in scientific
computing. Its ecosystem has become particularly important for numerical and
array-based computation, notably through libraries such as NumPy
\cite{harris2020array}. Although Python supports multiple programming styles,
numerical algorithms can naturally be expressed using conventional imperative
constructs such as loops, assignments, and conditional statements.

Haskell, in contrast, is a purely functional programming language characterized
by immutable values, first-class functions, higher-order functions, recursion,
strong static typing, and non-strict evaluation
\cite{marlow2010haskell,hutton2016programming}. These characteristics make
Haskell especially suitable for investigating how numerical algorithms can be
expressed without relying on explicit mutable state.

The use of these three languages allows different representations of the same
numerical algorithm to be investigated while exposing distinct computational
models.

Among the numerical algorithms considered in this work are iterative methods
such as fixed-point iteration and the Newton--Raphson method. Such algorithms
are particularly suitable for comparisons between programming paradigms
because their mathematical formulations are based on recursively defined
sequences of approximations
\cite{burden2016numerical,chapra2021numerical}.

In fixed-point iteration, for example, the objective is to determine a value
$x^\ast$ satisfying

\begin{equation}
    x^\ast = g(x^\ast).
\end{equation}

Starting from an initial approximation $x_0$, a sequence is generated according
to

\begin{equation}
    x_{n+1} = g(x_n),
\end{equation}

until a prescribed convergence criterion is satisfied.

The convergence of fixed-point iteration depends on mathematical properties of
the function $g$ and on the initial approximation. Under appropriate
conditions, the generated sequence converges to a fixed point of the function
\cite{burden2016numerical}.

This recursive structure highlights one of the central differences between the
programming paradigms investigated in this work. In an imperative
implementation, the current approximation may be explicitly updated during
each loop iteration. In a functional implementation, the sequence may instead
be described as successive applications of a function without requiring an
explicitly mutable variable.

The Newton--Raphson method provides another important example. Given a nonlinear
equation

\begin{equation}
    f(x)=0,
\end{equation}

an iterative approximation is obtained using

\begin{equation}
    x_{n+1}
    =
    x_n -
    \frac{f(x_n)}{f'(x_n)}.
\end{equation}

Under appropriate regularity assumptions and for sufficiently suitable initial
approximations, Newton's method can exhibit rapid local convergence
\cite{burden2016numerical,chapra2021numerical}.

Despite the mathematical equivalence of implementations based on the same
recurrence relations, their computational characteristics may differ. This
work does not attempt to establish one programming paradigm as universally
superior to another. Instead, it investigates quantitatively and qualitatively
how different implementation models behave in the particular context of
numerical algorithms.

The quantitative analysis considers metrics such as execution time, memory
usage, computational resource consumption, number of iterations, and
implementation size. Whenever applicable, compiler optimization levels and
language execution models are also taken into account.

Experimental comparisons between programming languages require particular
care. Differences in execution time cannot necessarily be attributed to the
programming paradigm alone, since performance is also influenced by compiler
technology, runtime implementation, optimization strategies, memory
management, data representation, and hardware architecture.

For this reason, the experiments presented in this work are conducted under
controlled conditions, using the same hardware platform and equivalent testing
conditions whenever possible. Different test functions and computational
scenarios are employed to examine the behavior of the implementations under
varying levels of numerical and computational complexity.

In addition to quantitative metrics, qualitative characteristics are examined,
including readability, modularity, expressiveness, maintainability, and the
degree of correspondence between the mathematical formulation of an algorithm
and its source-code representation.

Functional programming is particularly interesting in this context because
higher-order functions and abstraction mechanisms can improve program
modularity and permit mathematical operations to be represented compositionally
\cite{hughes1989functional}. However, properties such as recursion and lazy
evaluation may also result in execution and memory behaviors that differ from
those encountered in conventional imperative implementations.

A further component of this research concerns floating-point arithmetic.
Numerical algorithms are executed using finite computational representations,
which means that mathematical operations over real numbers are replaced by
operations over a finite subset of representable values. The consequences of
this approximation include rounding error, loss of significance, cancellation,
underflow, overflow, and other finite-precision effects
\cite{higham2002accuracy,ieee7542019}.

These effects are fundamental when interpreting the experimental results.
Differences between implementations must therefore not be analyzed solely in
terms of execution speed, but also with respect to the numerical behavior of
the algorithms and the arithmetic used by the underlying computational system.

The general objective of this work is to analyze the relationship between
programming paradigms and numerical methods, evaluating how different
implementation strategies influence both computational and numerical
properties of the algorithms.

The specific objectives are:

\begin{itemize}
    \item to study the mathematical foundations of the numerical methods
    considered in this work;

    \item to implement numerical algorithms using C, Python, and Haskell;

    \item to represent equivalent numerical methods using imperative and
    functional programming approaches;

    \item to compare the implementations in terms of execution time, memory
    consumption, and computational resource usage;

    \item to analyze differences in readability, modularity, and expressiveness
    among the implementations;

    \item to investigate the relationship between recursion, iterative control
    structures, and computational performance;

    \item to study the effects of floating-point representation on numerical
    algorithms;

    \item to investigate the propagation of representation and rounding errors
    throughout iterative numerical processes;

    \item to identify advantages and limitations associated with imperative and
    functional programming approaches in the implementation of numerical
    methods.
\end{itemize}

By combining these perspectives, this research aims to provide an analysis
that extends beyond a direct comparison of programming languages. The broader
objective is to understand how different computational models interact with
the mathematical characteristics of numerical algorithms and how those
interactions affect performance, implementation structure, and numerical
reliability.